\chapter{Спектральное кручение как последний селектор заряженных рёбер}

Правильно типизированный модуль одноформ сузил неоднозначность до двух
относительных координат между рёбрами $u,d,e$. Проверим последний внутренний
механизм, который не является новым потенциалом: способно ли спектральное
кручение или соответствующее исчисление второй степени не только видеть,
но и однозначно выбирать эти координаты.

\section{Антикруговая сверка}

Обычный фактор двухформ на ранг-один семейном углу уже вычислялся и оставил
лишь скалярную опорную линию. Повторять его нельзя. Новизна настоящего теста
состоит в использовании трёхлинейного спектрального функционала кручения и
модифицированного исчисления второй степени, предложенного для внутренней
геометрии Стандартной модели в \path{arXiv:2511.08159}; исходное определение
спектрального кручения дано в \path{arXiv:2308.01644}.

Литературный результат следует читать осторожно. Для уже заданного
конечного оператора Дирака спектральное кручение ненулевое. Обычное
исчисление Конна не позволяет буквально воспроизвести его связностью.
Совпадение достигается после замены произведения второй степени специальным
идемпотентом формализма Месланда--Ренни. Этот идемпотент не эрмитов и не
согласован с инволюцией; после дополнительного проектирования авторы
получают единственные алгебраические данные, воспроизводящие уже известный
спектральный функционал.

Следовательно, необходимо различать два утверждения:
\begin{enumerate}
 \item кручение восстанавливает функционал заданного $D_F$;
 \item кручение выводит неизвестный $D_F$.
\end{enumerate}
Второе из первого не следует.

\section{Ограничение на пакет \texorpdfstring{$H_{15}$}{H15}}

Для одного поколения обозначим
\begin{equation}
 a=|y_u|^2,\qquad b=|y_d|^2,\qquad c=|y_e|^2.
\end{equation}
После удаления правого нейтрино спектральное кручение содержит, в частности,
три независимых квартичных сочетания
\begin{equation}
 I_u=3a^2,\qquad
 I_{de}=3b^2+c^2,\qquad
 I_{ud}=3ab.
 \label{eq:v5-h15-torsion-invariants}
\end{equation}
Последнее слагаемое является настоящей межрёберной чувствительностью:
кручение не распадается на три полностью независимых одночастичных числа.

Однако все три величины однородны:
\begin{equation}
 (y_u,y_d,y_e)\longmapsto
 t(y_u,y_d,y_e)
 \quad\Longrightarrow\quad
 I\longmapsto t^4I.
\end{equation}
Поэтому сам тензор кручения не задаёт абсолютной амплитуды.

\section{Точный ранг относительной чувствительности}

После удаления общего масштаба положим
\begin{equation}
 r=\frac{a}{b},\qquad s=\frac{c}{b}.
\end{equation}
Две безмасштабные комбинации можно выбрать как
\begin{equation}
 R_1=\frac{I_{ud}}{I_u}=\frac1r,
 \qquad
 R_2=\frac{I_{de}}{3b^2}=1+\frac{s^2}{3}.
\end{equation}
Их якобиан равен
\begin{equation}
 \frac{\partial(R_1,R_2)}{\partial(r,s)}
 =
 \begin{pmatrix}
 -r^{-2}&0\\
 0&2s/3
 \end{pmatrix},
 \qquad
 \det=-\frac{2s}{3r^2}.
 \label{eq:v5-h15-torsion-relative-jacobian}
\end{equation}
На общем ненулевом фоне ранг равен двум. Значит, спектральное кручение
действительно различает обе оставшиеся относительные координаты. Но оно
возвращает их значения как функции уже заданных $y_u,y_d,y_e$ и не содержит
уравнения, которое предписывает конкретную точку $(r,s)$.

\begin{proposition}[Чувствительность не является выбором]
Полноранговое отображение от юкавских отношений к компонентам кручения
позволяет восстановить отношения из заранее измеренного тензора. Оно не
является вариационным или алгебраическим условием на эти отношения, пока
значения компонент кручения не заданы независимо.
\end{proposition}

\section{Почему положительная норма не спасает ветвь}

Можно формально записать
\begin{equation}
 V_T=I_u^2+I_{de}^2+I_{ud}^2\geq0.
\end{equation}
Но это уже новый потенциал. Без независимо выведенного ненулевого уровня
его глобальный минимум достигается при
\begin{equation}
 a=b=c=0.
\end{equation}
Добавление условия $I=I_0$, линейного источника или относительных весов
возвращает ту же внешнюю информацию, которую гейт должен был вывести.

\begin{theorem}[Запрет спектрального кручения как юкавского селектора]
На пакете $H_{15}$ спектральное кручение ненулево и имеет полный ранг два
по двум относительным заряженным юкавским координатам. Тем самым оно их
измеряет, но не выбирает. Алгебраическая связность, воспроизводящая
спектральный функционал в модифицированном исчислении, строится после
задания оператора Дирака и не выводит его коэффициенты. Обычное
инволютивное исчисление второй степени совпадения не даёт. Поэтому
спектральное кручение не замыкает юкавскую динамику родителя Мориты.
\end{theorem}

\begin{nogocriterion}[Запрет нового потенциала кручения]
Нельзя продолжать ветвь минимизацией произвольной нормы компонент кручения
или назначением им ненулевого уровня. Такой шаг вводит новый функционал или
новые наблюдаемые данные после обнаружения двумерной свободы.
\end{nogocriterion}

\begin{protocol}
\begin{enumerate}
 \item ненулевое кручение на $H_{15}$: \textbf{получено};
 \item чувствительность к смешению $u$--$d$: \textbf{получена};
 \item ранг по двум относительным координатам: \textbf{два};
 \item независимые значения компонент кручения:
 \textbf{не выведены};
 \item обычное инволютивное исчисление как точное воспроизведение:
 \textbf{не проходит};
 \item модифицированное исчисление как независимый селектор:
 \textbf{не проходит, поскольку использует заданный $D_F$};
 \item путь Мориты к единственному оператору Юкавы:
 \textbf{закрыт};
 \item родитель Мориты как кинематический контейнер:
 \textbf{сохранён};
 \item физическое замыкание: \textbf{не достигнуто}.
\end{enumerate}
\end{protocol}

Следующий шаг не должен вводить ещё один функционал на плоскости $(r,s)$.
Для нового открытия требуется независимо выведенная симметрия, связывающая
неэквивалентные рёбра, либо иной родительский принцип. До его появления
одноформенная ветвь Мориты замораживается как завершённый отрицательный
результат.

Машинный сертификат сохранён в
\path{s2t_v5_h15_spectral_torsion_selector_gate.py} и
\path{s2t_v5_h15_spectral_torsion_selector_gate_results.json}.